% Guideline & Policy for Project Proposal
% Only one project proposal needs to be submitted per group. The idea behind this proposal is
% that I can give you feedback and suggestions to make you come up with a feasible plan for the
% project. You don’t have to do exactly what you say in your proposal and can even completely
% change your solving strategies afterwards if you want. But it’s important to have at least one
% reasonable plan/baseline to start from. The length should be 4 to 6 pages, not including
% appendices.


% Structure
% It should contain title, authors, abstract, introduction, related work, model/method, and
% experiments. In particular, all these sections should be complete (though they can be short)
% except for model/method and experiments.
\documentclass{article}

% Use final option for non-anonymous class submission
\PassOptionsToPackage{numbers}{natbib}  % set numeric citations
\usepackage[final]{neurips_2021}

% Recommended packages
\usepackage[utf8]{inputenc}
\usepackage[T1]{fontenc}
\usepackage{hyperref}
\usepackage{url}
\usepackage{booktabs}
\usepackage{amsfonts}
\usepackage{nicefrac}
\usepackage{microtype}
\usepackage{xcolor}
\usepackage{amsmath}
\usepackage{graphicx}
\usepackage{longtable}

\title{Spring 2026 Analytics II Team Project Final Report}

% List the group and authors
\author{
  Group 1 \\[1ex]  % Adds Group 1 at the top
  Joshua Boyles\\
  Dave Julianto\\
  Hayden Burke\\
  Daniel Veiner\\
  Sebastian Hajdu
}

% Set the date manually 
%\date{May 1, 2026}

%\author{
%  First Author \\
%  Department of XYZ \\
%  University Name \\
%  City, Country \\
%  \texttt{author1@university.edu}
%  \And
%  Second Author \\
%  Department of XYZ \\
%  University Name \\
%  City, Country \\
%  \texttt{author2@university.edu}
%}


\begin{document}

\maketitle

% Add professor's name below title page (separate from authors)
\vspace{2em}  % Adds vertical space between the authors and professor's name
\begin{center}
  \textbf{Professor:}\\[0.5em] Dr. Kaixun Hua
\end{center}

\vspace{2em}
\begin{center}
  {\textbf{Date:}}\\[0.5em]
  May 1, 2026
\end{center}

% Abstract (15%)
% It should summarize the main idea of the project and its contributions. It should be
% understandable to anyone in the course.
\newpage
\begin{abstract}
%\section{Abstract}
Manufacturing defect detection plays a crucial role in maintaining product
quality and reliability in industrial systems. In ball bearing manufacturing,
defective components can lead to the following if not detected early in the
production process:

\begin{enumerate}
    \item equipment failures
    \item increased maintenance
    \item potential safety risks
\end{enumerate}

This project investigates the use of machine learning techniques
to predict defective ball bearings using historical manufacturing data
obtained from destructive and non-destructive testing procedures.

The dataset contains approximately 229,788 observations with 30 numerical
predictor variables representing sensor readings or process measurements,
along with a binary classification label indicating whether a bearing is
defective or non-defective. A major challenge associated with this dataset
is the severe class imbalance, as only 458 observations---approximately
0.2\% of the dataset---are labeled as defective. In such cases, traditional
metrics such as overall accuracy can be misleading because models may
achieve high accuracy by predicting the majority class while failing to
detect rare defects.

To address this issue, this project focuses on evaluation metrics that
better reflect performance on the minority class. Specifically, recall,
precision, F1-minority score, and the Area Under the Precision--Recall
Curve (PR-AUC) will be used to evaluate model performance, since missing
a defective component is significantly more costly than incorrectly
flagging a non-defective one.

Several techniques will be explored to mitigate class imbalance during
model training, including random oversampling, random undersampling,
Synthetic Minority Oversampling Technique (SMOTE), and class-weighted
learning methods. Multiple machine learning models will be evaluated,
including Logistic Regression, Decision Trees, Random Forests, and
boosting methods such as AdaBoost or gradient-boosted trees.

The experimental design includes model benchmarking, ablation studies of
imbalance-handling strategies, and sensitivity analysis across different
data split ratios. In addition, interpretability methods such as feature
importance analysis and SHAP values will be used to identify which
predictors contribute most strongly to defect detection.

The goal of this project is to develop a reliable and interpretable
machine learning framework that improves defect detection while
providing actionable insights for manufacturing quality control.
\end{abstract}


% Introduction (15%)
% It should clearly state the problem being addressed and why it is important.
% From ChatGPT: Introduce the research problem, motivation, and contributions of the paper.
\section{Introduction}

This project aims to utilize historical ball bearing manufacturing data to predict 
defectiveness of manufactured parts and to understand which combinations of control 
settings typically lead to defective production. The goal is to create an interpretable 
model that accurately predicts whether a ball bearing will be classified as defective 
after several destructive and non-destructive tests, based on a set of predictor variables.

The dataset contains approximately 230,000 samples with 30 predictor variables. Each sample 
is labeled as defective (class 1) or non-defective (class 0), which is the target variable 
we aim to predict. By understanding the factors that contribute to defects, corrective 
action can be taken on controllable system inputs to reduce future defects. Reduced defect 
rates directly lower scrap, rework, resource, and labor costs in the manufacturing process.

However, this problem presents several challenges: the dataset is large and high-dimensional, 
the classes are severely imbalanced (only around 0.2\% of samples defective), and manufacturing 
processes can involve complex, nonlinear interactions between variables. Additionally, 
interpretability is critical, as predictions must provide actionable insights for engineers 
and management.

In this project, we aim to:

\begin{enumerate}
    \item Develop an interpretable predictive model for classifying defective ball bearings.
    \item Identify and analyze key process variables that most strongly influence defect rates.
    \item Explore how the model’s insights can guide corrective actions to reduce defects and improve manufacturing efficiency.
\end{enumerate}

Together, these aims outline a plan to leverage historical manufacturing data for 
actionable defect prediction and process improvement.


% Related Work (15%)
% It should clearly distinguish what your proposed contribution from previous literature (e.g., a
% 1-2 sentence summary of at least 3 closely related papers).
\section{Related Work}
Machine learning has increasingly been applied to manufacturing systems for predictive quality monitoring and defect detection. Prior research shows that industrial process and test data can be used to learn relationships between operating conditions and product quality outcomes.

Zahrani et al.\ apply machine learning models to predict product quality and groove geometry in a circular laser-grooving manufacturing process. Their results show that data-driven models can capture nonlinear relationships between process parameters and final part characteristics, supporting predictive monitoring rather than purely reactive inspection \cite{zahrani2020lasergrooving}.

More recent studies extend this direction using stronger ensemble methods for rare-defect prediction in industrial environments. Geng et al.\ propose a framework for predicting internal cracks in continuous steel casting using CatBoost combined with synthetic data generation, demonstrating improved detection of infrequent defect patterns in a real manufacturing context \cite{geng2025ctgan}.

A central challenge in manufacturing defect prediction is severe class imbalance. Defective products may represent only a small fraction of observations, causing standard training and evaluation to favor the majority class and report high accuracy while missing rare failures. Kafunah et al.\ examine this issue in fault detection and evaluate imbalance-aware strategies such as oversampling and reweighting to improve minority-class detection, highlighting the importance of focusing on defect-class metrics such as precision, recall, and F1 rather than accuracy alone \cite{kafunah2021imbalanced}.

Cost-sensitive learning provides a complementary approach by explicitly assigning higher penalties to misclassifying defective samples, reflecting the higher operational risk of false negatives in manufacturing. Liu et al.\ introduce a cost-sensitive defect detection framework that modifies the training objective using an explicit cost structure to prioritize costly minority-class errors, and they report improved defect identification under imbalance \cite{liu2023costsensitiveyolov5}.

Building on these findings, our project systematically compares interpretable and ensemble models, including logistic regression, decision trees, random forests, and boosting under multiple imbalance-mitigation strategies. We emphasize cost-sensitive random forests and minority-class performance metrics to evaluate model effectiveness for defect detection in a highly imbalanced manufacturing dataset.



\section{Model and Method}
 
\subsection*{Class Imbalance - Preferred Metrics}
The dataset comprises 229,788 ball-bearing test records featuring 30 numerical predictors and a binary response variable. With only 458 defective samples, standard classifiers tend to favor the majority class; yielding high accuracy that can be misleading. In manufacturing, false negatives are more costly than false positives, as they can lead to diminished customer confidence, expensive product recalls, or liability lawsuits. Consequently, accuracy is not an appropriate metric. Instead, recall, precision, F1 Score, and the Area Under the Precision-Recall Curve (PR-AUC) are the primary metrics.
 
Because the defective class represents only 0.2\% of the dataset, using Receiver Operating Characteristic Area Under Curve (ROC AUC) can overly estimate the performance. A study was conducted and proved that ROC AUC is not dependable if the true positive sample is less than 3\% of the total \cite{imani2026roc}. PR AUC is preferred because it summarizes the precision and recall tradeoff across thresholds and reflects performance exclusively on the defective class, suitable for the case of missing a defect is far more consequential than flagging a non defective part. However, ROC AUC could still be considered to gauge the performance of baseline models.
 
Using regular F1 score as a metric combines performance across all classes, which causes the majority class to dominate the result in imbalanced datasets, shown by this formula:
$$F1_{regular} = 2 \cdot \frac{P_{overall} \cdot R_{overall}}{P_{overall} + R_{overall}}$$
where $P_{overall}$ and $R_{overall}$ are precision and recall across both classes. Because these aggregated values are driven primarily by the majority class, a model can achieve a high regular F1 score even while performing poorly on the defective class.
F1 per class, especially F1 minority avoids this masking effect by computing precision and recall only for the defective class:
$$F1_{minority} = 2 \cdot \frac{P_{minority} \cdot R_{minority}}{P_{minority} + R_{minority}}$$
The criticality of using F1 to detect the minority class in classification serves to aid in anomaly detection and identify the imbalance of the data set and its analysis \cite{sciencedirect_minority}. Classification trained on imbalanced data is weak at predicting samples within the minority class.
 
\subsection*{Class Imbalance - Modification on Training Data}
To address the imbalance, four manipulation strategies are employed, each with its distinct trade-offs. Random Oversampling duplicates minority samples to improve recall, though it risks overfitting. Random Undersampling removes majority samples to balance the classes but may lead to information loss. Synthetic Minority Oversampling Technique (SMOTE) generates synthetic defective samples by interpolation to refine decision boundaries but may introduce noise if the minority class has high variance. Finally, Class-Weighted training adjusts the learning objective to penalize minority-class errors more heavily, shifting the decision boundary without altering the dataset. These methods are systematically evaluated by comparing performance metrics across the experimental design combinations.
 
\subsection*{Model Considerations}
We plan to evaluate four primary models, which are logistic regression, Decision Tree, Random Forest, and AdaBoost/XGBoost. By comparing these simple and complex models, performance trade-offs between interpretability and prediction results can be measured.
 
\subsubsection*{Logistic Regression}
Logistic Regression is a sensible model to explore as it is suitable for binary classification problems. Logistic Regression models the probability that a bearing is defective as a smooth function of the readings:
$$P(defective|x) = \sigma(\beta_{0} + \beta^{\top}x) = \frac{1}{1 + e^{-(\beta_{0} + \beta^{\top}x)}}$$
Here, $x$ is the vector of reading values, $\beta$ is the vector of coefficients, and $\sigma(\cdot)$ is the logistic (sigmoid) function that maps any real number to a probability between 0 and 1. Each coefficient $\beta_i$ represents the change in the log-odds of a defect associated with a one-unit increase in reading $x_i$, holding other predictors constant. This makes Logistic Regression highly interpretable as engineers can directly see which predictors increase or decrease defect risk and by how much.
 
\subsubsection*{Decision Trees}
Decision Trees provide transparency by sorting data through a series of logical, binary splits. At each step, the model identifies the sensor threshold that best separates defective bearings from healthy ones, reducing impurity until a homogeneous classification is reached. This visual approach allows us to see exactly which predictor values trigger a defect alert. To determine the best possible split, the algorithm calculates a purity score at each step. This is done using Gini Impurity, which measures the probability of a specific sensor reading being incorrectly classified if it were labeled randomly. A lower Gini value indicates a more pure or successful separation of defective bearings from healthy ones. Gini is calculated as follows:
$$Gini(S) = \sum_{k \in \{0,1\}} p_{k}(1 - p_{k})$$
where $p_k$ is the proportion of class $k$ in node $S$. A split into left and right child nodes $S_L$ and $S_R$ is evaluated by:
$$\Delta = Gini(S) - \left(\frac{|S_{L}|}{|S|}Gini(S_{L}) + \frac{|S_{R}|}{|S|}Gini(S_{R})\right)$$
The algorithm selects the split with the largest $\Delta$, the greatest reduction in impurity.
Gini impurity is used instead of entropy because it is computationally simpler by not
using logarithms and tends to be more numerically stable under extreme class imbalance.
As the defective class is severely smaller, tiny fluctuations in estimated probabilities
can cause large changes in entropy, potentially amplifying noise. Gini provides smoother
behavior, leading to more stable splits and more reliable rankings of feature importance.
For this project, Decision Tree models serve as an interpretable nonlinear model that
can capture simple interactions between predictors, and as a base learner for ensemble methods.
 
\subsubsection*{Random Forest}
Random Forest extends Decision Trees by aggregating many trees trained on different subsets of the data and features. Each tree is trained on a bootstrap sample of the training data, and at each split, a random subset of features is considered. The final prediction is obtained by majority vote:
$$\hat{y} = mode\{\hat{y}^{(1)}, \hat{y}^{(2)}, \dots, \hat{y}^{(N)}\}$$
where $\hat{y}^{(n)}$ is the prediction of tree $n$, and $N$ is the number of trees. The ensemble model determines the best prediction by calculating the mode across the results of each individual tree. To address class imbalance, the Balanced Random Forest variant is used: for each tree, all defective samples are included, and an equal number of non-defective samples are randomly drawn to form a balanced bootstrap sample. This ensures that each tree is trained on a dataset where the defective class is not overwhelmed by the majority class, improving the model's ability to learn splits that separate defects from non-defects and thereby improving recall. Random Forests also provide feature importance scores (e.g., based on impurity reduction or permutation importance), which are valuable for manufacturing engineers who need to understand which sensors consistently contribute to defect detection across many trees.
 
While Balanced RF handles imbalance by resampling the data for each tree, cost-sensitive RF addresses the same issue by applying class weighted impurity so that errors on defective samples receive higher penalties (One application of the Class weighted strategy). Evaluating both RF variants allows us to compare the best model for detection performance.
 
\subsubsection*{Boosting}
Adaptive Boosting (AdaBoost) builds an ensemble of weak learners by repeatedly reweighting the training samples so that each new classifier focuses on the mistakes of the previous ones. The misclassified samples get their weights $w_j^i$ increased, correctly classified ones decreased, and then they all get normalized by the normalization constant. Over time, this forces the ensemble to concentrate on high-error instances:
$$\epsilon_i = \sum_{j} w_j^i \delta(C_i(x_j) \neq y_j)$$
This quantity measures how much weight the current classifier assigns to its mistakes. The classifier’s importance weight is then computed as:
$$\alpha_{t} = \frac{1}{2} \ln\left(\frac{1 - \epsilon_{t}}{\epsilon_{t}}\right)$$
The importance increases when the classifier performs better than chance and decreases when it performs poorly. After each iteration, AdaBoost updates the sample weights:
$$w_{j}^{i+1} = \frac{w_{j}^{i}}{Z_{i}} \times \begin{cases} e^{-\alpha_{i}}, & C_{i}(x_{j}) = y_{j} \\ e^{\alpha_{i}}, & C_{i}(x_{j}) \neq y_{j} \end{cases}$$
The misclassified samples get their weights increased by $e^{\alpha_i}$, correctly classified ones decreased by $e^{-\alpha_i}$, and then they all get normalized by $Z_i$, the normalization constant. Over time, this forces the ensemble to concentrate on high error samples. This behavior is useful for classifying rare defective bearings in our data. As they are rare, they are often misclassified by early weak learners because they appear infrequently and have subtle patterns. AdaBoost naturally increases their weight, causing them to gain importance. This can improve defect detection without manually oversampling. However, AdaBoost is sensitive to noise; overemphasizing and overfitting if there is too much. For this reason, AdaBoost serves as a strong but potentially sensitive benchmark in the project timeline.
 
AdaBoost will serve as our baseline boosting methods, providing a simple and interpretable reference for boosting model. From here, we plan to extend the analysis to more advanced gradient boosting models such as XGBoost, LightGBM, and CatBoost to gauge performance gains, interpretability, and process time. All boosting variants will be included in the model benchmark, with XGBoost being the prioritized model to be explored.
 

\section{Experimental Design}
 
The experimental phase is structured into three distinct components to ensure a robust evaluation of the model.
 
\begin{itemize}
    \item \textbf{Model Benchmark:} A comparison of simple and complex models to determine the best performing model; with consideration of interpretability vs. predictive performance.
   
    \item \textbf{Class Imbalance Strategy Ablation:} A systematic evaluation of various dataset manipulation strategies to address the class imbalance problem.
   
    \item \textbf{Data Split Comparison:} A test of model stability and generalization across different training, validation, and test ratios.
\end{itemize}
 
\subsection*{Main Benchmark Analysis}
 
\textbf{Objective:} Evaluate candidate models... identify the optimal balance between predictive power and model transparency... This benchmark compares the interpretability vs. performance of models... and is often as critical as its accuracy in industrial maintenance environments.
Something about focus on false negative rates and false positive rates?
 
\subsection*{Evaluation Focus:}
 
something about the scoring system...
"FN \& FP Score"
 
 
\subsubsection*{Expected Outcomes:}
 
\begin{itemize}
    \item \textbf{Superiority of Ensembles:} Boosting and Forest based methods are expected to yield the highest PR-AUC.
    \item \textbf{Diminishing Returns:} We anticipate a curve where complex boosting methods provide only marginal gains over Balanced Random Forests, but at the cost of significantly higher complexity and problems due to noisy minority samples.
    \item \textbf{Interpretability:} Baseline (Simple) models will likely provide the highest interpretability but lower recall, serving as the essential control group for the study.
\end{itemize}
 
\subsection*{Class Imbalance Strategy Ablation}
 
\textbf{Objective:} Quantify the impact of various sampling and weighting strategies on the decision boundary. Each of the strategies will be tested to each of the model selection, creating 16 experimental combinations.
 
\begin{itemize}
    \item \textbf{Oversampling} Check whether duplicating some minority samples boost recall without overfitting.
    \item \textbf{Undersampling:} Check whether removing some majority samples improves performance and interpretability balance without severely losing information.
    \item \textbf{SMOTE:} Testing if synthetic minority generation improves recall without introducing excessive noise.
    \item \textbf{Cost-Sensitive Learning:} Evaluating if class weighting offers a cleaner alternative to data manipulation.
    \begin{itemize}
    \item \textbf{Logistic Regression (Class-Weighted Loss)}: Applies class weights inside the weighted cross-entropy objective so that misclassifying defective samples incurs a higher penalty.
 
    \item \textbf{Decision Tree (Class-Weighted Impurity)}: Uses class-weighted Gini impurity, causing splits to favor isolating minority-class samples.
 
    \item \textbf{Random Forest (Cost-Sensitive RF)}: Implements class weighting through weighted impurity across all trees; this is the Random-Forest-specific form of class-weighted learning.
 
    \item \textbf{Boosting (AdaBoost/XGBoost) with Class Weights}: Adjusts sample weights or uses parameters such as \texttt{scale\_pos\_weight} (multiplying training loss to increase minority class influence).
\end{itemize}
 
\end{itemize}
 
\subsubsection*{Expected Outcome:}
 
Class weighting and Balanced Random Forest (internal undersampling) are expected to provide the most stable results, whereas SMOTE may increase False Positives by creating unrealistic defective signatures.
 
\subsubsection*{Data Split Comparison}
 
This test is to determine if the model’s performance is affected by a specific data split, especially given the limited minority sample size. The baseline data split will use 80/10/10 of training, validation, and test data, respectively.
 
Remember to add the other splits!!!!!

\section{Results}

\section*{Boosting Models Experiment}

The first complete run of the experiment was done with default parameters for all boosting models. A second run was made with the tuned hyperparameters of boosting models determined by a 5-Fold Cross Validation on the 80/20 baseline split resulting in:

\begin{table}[h!]
\centering
\caption{Boosting Models Hyperparameters: Default vs Tuned}
\begin{tabular}{l l l l}
\hline
\textbf{Model} & \textbf{Parameter} & \textbf{Default} & \textbf{Tuned} \\
\hline
AdaBoost & n\_estimators & 50 & 200 \\
         & learning\_rate & 1.0 & 1.5 \\
XGBoost  & n\_estimators & 100 & 300 \\
         & learning\_rate & 0.3 & 0.05 \\
LightGBM & n\_estimators & 100 & 300 \\
CatBoost & iterations & 1000 & 300 \\
\hline
\end{tabular}
\end{table}

\subsection{Ablation Analysis}

The data split ratio has only minor effects on aggregate performance. The class imbalance treatment is the overwhelming driver of performance, especially on the precision-recall tradeoff.

\paragraph{Random Undersampling:} There is a severe overfit / distribution mismatch. Nearly all four models hit train precision of exactly 1.0 but have less than 10\% test precision, with recall mostly above 70\%. This indicates overfitting as the model memorizes a balanced but artificially shrunk training distribution, then encounters the natural imbalance at test time and floods predictions toward the minority class. This may also indicate a fundamental mismatch between training and deployment distributions.

\paragraph{Baseline (Unweighted):} XGBoost and CatBoost both perform strongly. Train and test precision align closely (\~96–100\% test), recall sits around 70–80\%, and ROC AUC stays above 0.98 — a genuinely stable result with minimal overfit. AdaBoost behaves similarly but with lower recall (57–70\%), suggesting it is conservative rather than overfit. LightGBM is the outlier: performance degrades dramatically as the training set shrinks, going from near-collapse at 70/30 (test precision 10.8\%, recall 15.9\%, ROC AUC 0.58) to merely mediocre at 90/10. LightGBM without any balancing strategy appears sensitive to the combination of severe class imbalance and limited training data.

\paragraph{Class-Weighted:} AdaBoost does not respond. For XGBoost, CatBoost, and LightGBM, class-weighting produces a meaningful lift in recall (\~76–80\%) while keeping test precision above 92\%, making it a clean, low-complexity improvement over the unweighted baseline. However, AdaBoost's class-weighted results are bit-for-bit identical to its unweighted baseline across all three splits — scikit-learn's AdaBoostClassifier does not accept a \texttt{class\_weight} parameter, so the two configurations are effectively running the same model twice.

\paragraph{SMOTE:} XGBoost, CatBoost, and LightGBM all achieve a reasonable middle ground under SMOTE: train precision near 1.0 with test precision in the 60–94\% range and recall around 75–85\%. The train-test precision gap (roughly 30–40 points) signals moderate overfitting to synthetic samples, but F1 scores are still competitive. AdaBoost is an outlier: test precision collapses to \~9\% while runtime balloons to 255–449 seconds — over 100× slower than CatBoost or XGBoost under the same conditions.

\paragraph{Random Oversampling:} XGBoost, LightGBM, and CatBoost achieve the best overall F1 scores, with test precision of 93–100\% and recall of 73–84\%. Train precision is 1.0 for all, and test precision remains high, with a small distribution shift compared to SMOTE. AdaBoost again lags substantially (21–37\% test precision), showing poor response to any resampling that inflates training set size.



\section{Conclusion and Future Work }

\subsection{From Boosting Summary...}

The clearest takeaways are: 
\begin{enumerate}
    \item Sampler choice is what matters, not how the data is split.
    \item Random undersampling produces a consistent overfit / distribution-mismatch signature across all four models.
    \item XGBoost and CatBoost are robust across all strategies.
    \item LightGBM should not be used without a balancing strategy given its baseline instability.
    \item AdaBoost largely fails to benefit from class-weighting, degrades under SMOTE and oversampling, and is by far the slowest — making it the weakest option among the four boosters evaluated here.
\end{enumerate}


%Bibliography
\newpage

\bibliographystyle{plain}
\begin{thebibliography}{9}

\bibitem{zahrani2020lasergrooving}
E.~G. Zahrani, F.~Hojati, A.~Daneshi, B.~Azarhoushang, and J.~Wilde.
\newblock Application of machine learning to predict the product quality and geometry in circular laser grooving process.
\newblock {\em Procedia CIRP}, 94:474--480, 2020.

\bibitem{geng2025ctgan}
M.~Geng, H.~Ma, S.~Liu, Z.~Zhou, L.~Xing, Y.~Ai, and W.~Zhang.
\newblock A data-driven approach for internal crack prediction in continuous casting of HSLA steels using CTGAN and CatBoost.
\newblock {\em Materials}, 18:3599, 2025.

\bibitem{kafunah2021imbalanced}
J.~Kafunah, M.~I. Ali, and J.~G. Breslin.
\newblock Handling imbalanced datasets for robust deep neural network-based fault detection in manufacturing systems.
\newblock {\em Applied Sciences}, 11:9783, 2021.

\bibitem{liu2023costsensitiveyolov5}
B.~Liu, F.~Gao, and Y.~Li.
\newblock Cost-sensitive YOLOv5 for detecting surface defects of industrial products.
\newblock {\em Sensors}, 23:2610, 2023.

\bibitem{imani2026roc}
M.~Imani, M.~Joudaki, A.~Bagheri, and H.~R. Arabnia.
\newblock Why {ROC-AUC} is misleading for highly imbalanced data: In-depth
  evaluation of {MCC}, {F2-Score}, {H-Measure}, and {AUC}-based metrics across
  diverse classifiers.
\newblock {\em Technologies}, 14(1):54, 2026.

\bibitem{sciencedirect_minority}
ScienceDirect.
\newblock Minority class.
\newblock {\em ScienceDirect Topics}.
\newblock Retrieved March 5, 2025.

\end{thebibliography}



%Appendix
\newpage
\section*{Appendix}
\subsection*{A. Detailed Results}


\begin{center}
\begin{longtable}{c c p{5cm} c c c}
\caption{"FN \& FP Scores" and Error Rates for Various Scenarios} \\
\toprule
FN Weight & Rank & Scenario & "FN \& FP Score" & FNR & FPR \\
\midrule
\endfirsthead
\toprule
FN Weight & Rank & Scenario & "FN \& FP Score" & FNR & FPR \\
\midrule
\endhead
0.5 & 1 & 90/10 | Random Undersampling | CatBoost & 0.034292137 & 0.045454545 & 0.023129728 \\
0.5 & 2 & 90/10 | Random Oversampling | Logistic Regression & 0.044256621 & 0.068181818 & 0.020331424 \\
0.5 & 3 & 90/10 | Class-Weighted | Logistic Regression & 0.044322207 & 0.068181818 & 0.020462595 \\
0.5 & 4 & 80/20 | SMOTE | Logistic Regression & 0.045513672 & 0.068181818 & 0.022845525 \\
0.5 & 5 & 90/10 | SMOTE | Logistic Regression & 0.045524603 & 0.068181818 & 0.022867387 \\
0.5 & 6 & 90/10 | Random Undersampling | AdaBoost & 0.046884502 & 0.045454545 & 0.048314459 \\
0.5 & 7 & 90/10 | Random Undersampling | LightGBM & 0.047186095 & 0.068181818 & 0.026190372 \\
0.5 & 8 & 90/10 | Random Undersampling | XGBoost & 0.047448436 & 0.068181818 & 0.026715054 \\
0.5 & 9 & 80/20 | Random Undersampling | Logistic Regression & 0.047538119 & 0.056818182 & 0.038258056 \\
0.5 & 10 & 90/10 | Random Undersampling | Logistic Regression & 0.047907533 & 0.068181818 & 0.027633247 \\
0.6 & 1 & 90/10 | Random Undersampling | CatBoost & 0.036524618 & 0.045454545 & 0.023129728 \\
0.6 & 2 & 90/10 | Random Undersampling | AdaBoost & 0.046598511 & 0.045454545 & 0.048314459 \\
0.6 & 3 & 90/10 | Random Oversampling | Logistic Regression & 0.04904166 & 0.068181818 & 0.020331424 \\
0.6 & 4 & 90/10 | Class-Weighted | Logistic Regression & 0.049094129 & 0.068181818 & 0.020462595 \\
0.6 & 5 & 80/20 | Random Undersampling | Logistic Regression & 0.049394132 & 0.056818182 & 0.038258056 \\
0.6 & 6 & 80/20 | SMOTE | Logistic Regression & 0.050047301 & 0.068181818 & 0.022845525 \\
0.6 & 7 & 90/10 | SMOTE | Logistic Regression & 0.050056046 & 0.068181818 & 0.022867387 \\
0.6 & 8 & 90/10 | Random Undersampling | LightGBM & 0.05138524 & 0.068181818 & 0.026190372 \\
0.6 & 9 & 90/10 | Random Undersampling | XGBoost & 0.051595112 & 0.068181818 & 0.026715054 \\
0.6 & 10 & 90/10 | Random Undersampling | Logistic Regression & 0.05196239 & 0.068181818 & 0.027633247 \\
0.7 & 1 & 90/10 | Random Undersampling | CatBoost & 0.0387571 & 0.045454545 & 0.023129728 \\
0.7 & 2 & 90/10 | Random Undersampling | AdaBoost & 0.046312519 & 0.045454545 & 0.048314459 \\
0.7 & 3 & 80/20 | Random Undersampling | Logistic Regression & 0.051250144 & 0.056818182 & 0.038258056 \\
0.7 & 4 & 90/10 | Random Oversampling | Logistic Regression & 0.0538267 & 0.068181818 & 0.020331424 \\
0.7 & 5 & 90/10 | Class-Weighted | Logistic Regression & 0.053866051 & 0.068181818 & 0.020462595 \\
0.7 & 6 & 80/20 | SMOTE | Logistic Regression & 0.05458093 & 0.068181818 & 0.022845525 \\
0.7 & 7 & 90/10 | SMOTE | Logistic Regression & 0.054587489 & 0.068181818 & 0.022867387 \\
0.7 & 8 & 90/10 | Random Undersampling | LightGBM & 0.055584384 & 0.068181818 & 0.026190372 \\
0.7 & 9 & 90/10 | Random Undersampling | XGBoost & 0.055741789 & 0.068181818 & 0.026715054 \\
0.7 & 10 & 90/10 | Random Undersampling | Logistic Regression & 0.056017247 & 0.068181818 & 0.027633247 \\
0.8 & 1 & 90/10 | Random Undersampling | CatBoost & 0.040989582 & 0.045454545 & 0.023129728 \\
0.8 & 2 & 90/10 | Random Undersampling | AdaBoost & 0.046026528 & 0.045454545 & 0.048314459 \\
0.8 & 3 & 80/20 | Random Undersampling | Logistic Regression & 0.053106157 & 0.056818182 & 0.038258056 \\
0.8 & 4 & 90/10 | Random Oversampling | Logistic Regression & 0.058611739 & 0.068181818 & 0.020331424 \\
0.8 & 5 & 90/10 | Class-Weighted | Logistic Regression & 0.058637973 & 0.068181818 & 0.020462595 \\
0.8 & 6 & 80/20 | SMOTE | Logistic Regression & 0.059114559 & 0.068181818 & 0.022845525 \\
0.8 & 7 & 90/10 | SMOTE | Logistic Regression & 0.059118932 & 0.068181818 & 0.022867387 \\
0.8 & 8 & 90/10 | Random Undersampling | LightGBM & 0.059783529 & 0.068181818 & 0.026190372 \\
0.8 & 9 & 90/10 | Random Undersampling | XGBoost & 0.059888465 & 0.068181818 & 0.026715054 \\
0.8 & 10 & 90/10 | Random Undersampling | Logistic Regression & 0.060072104 & 0.068181818 & 0.027633247 \\
0.9 & 1 & 90/10 | Random Undersampling | CatBoost & 0.043222063 & 0.045454545 & 0.023129728 \\
0.9 & 2 & 90/10 | Random Undersampling | AdaBoost & 0.045740536 & 0.045454545 & 0.048314459 \\
0.9 & 3 & 80/20 | Random Undersampling | Logistic Regression & 0.054962169 & 0.056818182 & 0.038258056 \\
0.9 & 4 & 90/10 | Random Oversampling | Logistic Regression & 0.063396779 & 0.068181818 & 0.020331424 \\
0.9 & 5 & 90/10 | Class-Weighted | Logistic Regression & 0.063409896 & 0.068181818 & 0.020462595 \\
0.9 & 6 & 80/20 | SMOTE | Logistic Regression & 0.063648189 & 0.068181818 & 0.022845525 \\
0.9 & 7 & 90/10 | SMOTE | Logistic Regression & 0.063650375 & 0.068181818 & 0.022867387 \\
0.9 & 8 & 90/10 | Random Undersampling | LightGBM & 0.063982673 & 0.068181818 & 0.026190372 \\
0.9 & 9 & 90/10 | Random Undersampling | XGBoost & 0.064035142 & 0.068181818 & 0.026715054 \\
0.9 & 10 & 90/10 | Random Undersampling | Logistic Regression & 0.064126961 & 0.068181818 & 0.027633247 \\
1 & 1 & 90/10 | Random Undersampling | CatBoost & 0.045454545 & 0.045454545 & 0.023129728 \\
1 & 2 & 90/10 | Random Undersampling | AdaBoost & 0.045454545 & 0.045454545 & 0.048314459 \\
1 & 3 & 80/20 | Random Undersampling | Logistic Regression & 0.056818182 & 0.056818182 & 0.038258056 \\
1 & 4 & 80/20 | Random Undersampling | CatBoost & 0.068181818 & 0.068181818 & 0.027720694 \\
1 & 5 & 80/20 | Random Undersampling | LightGBM & 0.068181818 & 0.068181818 & 0.030693892 \\
1 & 6 & 90/10 | Random Oversampling | Logistic Regression & 0.068181818 & 0.068181818 & 0.020331424 \\
1 & 7 & 90/10 | Random Undersampling | LightGBM & 0.068181818 & 0.068181818 & 0.026190372 \\
1 & 8 & 90/10 | SMOTE | Logistic Regression & 0.068181818 & 0.068181818 & 0.022867387 \\
1 & 9 & 80/20 | SMOTE | Logistic Regression & 0.068181818 & 0.068181818 & 0.022845525 \\
1 & 10 & 90/10 | Class-Weighted | Logistic Regression & 0.068181818 & 0.068181818 & 0.020462595 \\
\bottomrule
\end{longtable}
\end{center}


\end{document}
